\section{Evola on Will}

In The Individual and the Becoming of the World, Evola devotes himself to the question of the Will. While acknowledging the influence of Nietzsche on his thought, Evola nevertheless rejects Nietzsche's naturalism and reintroduces the supernatural element. He does this by bringing back Schopenhauer's second foundation of the world, the ``Idea'', which Evola calls ``essence'' in line with Traditional thought.

Evola rejects, however, Schopenhauer's view that consciousness is the passive observer of phenomena, which is a creation of the Will. In place of this dualistic view, Evola identifies the center of consciousness, the ``I'', with the Will. Thus the individual is both the creator and the observer of the World. For the undeveloped man, this will is ``spontaneous'', that is, it is unconscious, undirected, and predestined, rather like Schopenhauer's Will. (See the chapter ``Privation''\footnote{\url{https://gornahoor.net/Individual/IABW.htm\#Privation}}.)

On the other hand, the developed man has True Will and is the unconditioned cause of his world. As Evola puts it:

\begin{quotex}
This principle is: THE POWER OF CONTROL. The ``I'', in fact, is not a thing, a ``given'', a ``fact'', but, essentially, a deep centre of will and power. As ``I'', Fichte says, it is, only insofar as it posits itself --- and only a pure self-positing is, to tell the truth, its ``being''.
\end{quotex}


Note, in particular, that the ``I'' is not a fact, i.e., an object in the world, but is transcendent to the world.

As for the ``essence'' or ``concept'', Evola points out that concepts are sufficient for rational explanations; that is, the essence tells us ``what'' something is. However, the brute fact of existence --- ``that'' something is --- is not, and remains beyond rational explanation. Yet the fact of existence does become comprehensible when we realize that it is the creation of our Will. Hence, the measure of one's will power is the ability to actualize what was only there in potential.

So, Evola's supernaturalism is world-affirming, not world-denying. He does not advocate a passive escape to the pure world of ideas, but rather the active bringing the ideas into reality. Thusly does Evola succeed in explaining Nietzsche's Will to Power, while retaining the transcendent supernaturalism of Traditional metaphysics.


\flrightit{Posted on 2009-08-13 by Cologero}

